\chapter{QEMU Virtualization}
\label{ch:qemu}

QEMU (Quick EMUlator) is the primary test and deployment platform for
Veilbox. This chapter documents the QEMU command-line options, virtual
hardware configuration, network setup, and debugging techniques used
during development.

\section{QEMU Overview}

QEMU is a type 2 hypervisor that emulates a complete x86\_64 computer
system. When used with KVM (Kernel-based Virtual Machine), it achieves
near-native performance by executing guest instructions directly on the
host CPU.

\subsection{Boot Modes}

Veilbox supports two QEMU boot modes:

\paragraph{Direct Kernel Boot (-kernel)}:
\begin{lstlisting}
qemu-system-x86_64 \
    -kernel output/vmlinuz \
    -append "console=ttyS0 console=tty1" \
    -m 2G \
    -nographic
\end{lstlisting}

In this mode, QEMU loads the kernel binary directly without involving
a bootloader. The kernel must have \texttt{CONFIG\_EFI\_STUB=y} to
provide a valid MZ/PE header that QEMU's firmware recognizes.

\paragraph{BIOS Boot (SeaBIOS + GRUB)}:
\begin{lstlisting}
qemu-system-x86_64 \
    -drive file=output/veilbox.raw,format=raw,if=virtio \
    -m 2G \
    -nographic
\end{lstlisting}

In this mode, QEMU boots SeaBIOS, which loads GRUB from the MBR of the
disk image, which then loads the kernel. This path validates the full
boot process including the bootloader.

\section{Virtual Hardware Configuration}

\subsection{CPU and Memory}

\begin{lstlisting}
-cpu host        # Expose host CPU features to guest
-smp 2           # 2 virtual CPUs
-m 2G            # 2 GB RAM (required for ~90MB initramfs)
\end{lstlisting}

The \texttt{-cpu host} flag enables KVM's maximum feature set (VMX, SMEP,
SMAP, AES-NI, etc.). Without it, QEMU emulates a generic Nehalem-era CPU,
which lacks features that container runtimes may require.

\subsection{Storage}

\begin{lstlisting}
-drive file=output/veilbox.raw,format=raw,if=virtio
\end{lstlisting}

The VirtIO block device provides paravirtualized I/O. The guest kernel
accesses it as \texttt{/dev/vda}. Compared to IDE emulation
(\texttt{if=ide}), VirtIO block reduces I/O latency by 50\% and increases
throughput by 3--5x.

\subsection{Networking}

\begin{lstlisting}
-netdev user,id=net0,hostfwd=tcp::2222-:22
-device virtio-net,netdev=net0
\end{lstlisting}

User-mode networking (SLiRP) provides:
\begin{itemize}[nosep]
  \item A virtual DHCP server at 10.0.2.2.
  \item DNS redirection to the host's DNS (at 10.0.2.3).
  \item Guest IP: 10.0.2.15/24.
  \item Port forwarding: host port 2222 $\rightarrow$ guest port 22
        (SSH).
\end{itemize}

\section{Serial Console Modes}

\subsection{-nographic Mode}

\begin{lstlisting}
-nographic
\end{lstlisting}

QEMU redirects the serial port (ttyS0) to the terminal's stdio. The
guest output appears directly in the terminal. Keystrokes are sent to
the guest. This is headless mode --- no graphical window appears.

\subsection{Graphical Mode}

\begin{lstlisting}
-vga std        # Standard VGA adapter
-serial mon:telnet::2223,server,nowait  # Serial via telnet
\end{lstlisting}

With a graphical window, the guest's VGA text console (tty1) appears in
the QEMU window. The serial console is accessible via telnet on port 2223.

\section{QEMU Monitor}

QEMU provides a monitor interface for VM introspection:

\begin{lstlisting}
-monitor telnet::2224,server,nowait
\end{lstlisting}

Useful monitor commands:
\begin{lstlisting}
(qemu) info block        # List block devices and stats
(qemu) info network      # Network device statistics
(qemu) info registers    # CPU register dump
(qemu) info cpus         # CPU state
(qemu) system_powerdown  # Graceful shutdown
(qemu) quit              # Force quit
\end{lstlisting}

In \texttt{-nographic} mode, the monitor is accessed via \texttt{Ctrl-A C}.

\section{Test Script}

The \texttt{test.sh} script wraps the QEMU invocation with convenience
options:

\begin{lstlisting}[language=bash]
#!/bin/bash
# Direct kernel boot mode
MEM=${MEM:-2G}

qemu-system-x86_64 \
    -kernel output/vmlinuz \
    -append "console=ttyS0 console=tty1" \
    -m "$MEM" \
    -nographic \
    -netdev user,id=net0,hostfwd=tcp::2222-:22 \
    -device virtio-net,netdev=net0 \
    -drive file=output/state.img,format=raw,if=virtio
\end{lstlisting}

\section{Performance Tuning}

\subsection{KVM Acceleration}

Ensure the user is in the \texttt{kvm} group:
\begin{lstlisting}[language=bash]
sudo usermod -a -G kvm $USER
# Re-login or run: newgrp kvm
\end{lstlisting}

\subsection{Huge Pages}

Transparent Huge Pages (THP) improve memory performance:
\begin{lstlisting}[language=bash]
echo always > /sys/kernel/mm/transparent_hugepage/enabled
\end{lstlisting}

\subsection{CPU Pinning}

For dedicated VMs, pin vCPUs to physical cores:
\begin{lstlisting}
-smp 2,cores=2,threads=1 \
-object memory-backend-file,id=mem,size=2G,mem-path=/dev/hugepages \
-numa node,memdev=mem
\end{lstlisting}

\section{Advanced QEMU Options}

\subsection{Multiple Network Interfaces}

Configure multiple NICs for advanced container networking:

\begin{lstlisting}
-netdev user,id=net0,hostfwd=tcp::2222-:22 \
-netdev tap,id=net1,ifname=tap0,script=no,downscript=no \
-device virtio-net,netdev=net0,mac=52:54:00:12:34:56 \
-device virtio-net,netdev=net1,mac=52:54:00:12:34:57
\end{lstlisting}

\subsection{USB Passthrough}

For bare-metal testing, USB devices can be passed to the guest:

\begin{lstlisting}
-usb -device usb-host,vendorid=0x1234,productid=0x5678
\end{lstlisting}

\subsection{VM Snapshots}

QEMU supports VM snapshots for quick save/restore:

\begin{lstlisting}
# Start with snapshot support
qemu-system-x86_64 ... -drive file=veilbox.raw,format=raw,snapshot=on

# Create snapshot (in monitor)
(qemu) savevm my_snapshot

# Restore snapshot
(qemu) loadvm my_snapshot
\end{lstlisting}

\subsection{Verbose Debug Output}

Enable maximum debugging:

\begin{lstlisting}
-d cpu_reset,int,guest_errors,unimp \
-trace events=/tmp/trace-events
\end{lstlisting}

\section{Disk Image Management}

\subsection{Resizing the Disk Image}

To increase the disk image size:

\begin{lstlisting}[language=bash]
# Resize raw image to 16 GB
qemu-img resize output/veilbox.raw 16G

# Resize the partition
echo ', +' | sfdisk -N 1 output/veilbox.raw

# Resize the filesystem
# This requires loop mount or debugfs resize
\end{lstlisting}

\subsection{Converting Image Formats}

\begin{lstlisting}[language=bash]
# Raw to VDI (VirtualBox)
qemu-img convert -f raw -O vdi veilbox.raw veilbox.vdi

# Raw to VMDK (VMware)
qemu-img convert -f raw -O vmdk veilbox.raw veilbox.vmdk

# Raw to QCOW2 (QEMU native, compressed)
qemu-img convert -f raw -O qcow2 veilbox.raw veilbox.qcow2

# Raw to VHD (Hyper-V)
qemu-img convert -f raw -O vpc veilbox.raw veilbox.vhd

# Raw to compressed raw
qemu-img convert -f raw -O raw -c veilbox.raw veilbox-compressed.raw
\end{lstlisting}

\section{Health Check Mode}

The \texttt{--check} mode in test.sh automatically verifies that the
system boots correctly. It uses \texttt{expect} to detect the login
prompt:

\begin{lstlisting}[language=bash]
#!/bin/bash
# --check mode: auto-detect login prompt
expect << EOF
set timeout 30
spawn qemu-system-x86_64 -kernel output/vmlinuz -append \
    "console=ttyS0 console=tty1" -m 2G -nographic \
    -netdev user,id=net0 -device virtio-net,netdev=net0
expect {
    "veilbox login:" { exit 0 }
    timeout { exit 1 }
}
EOF
\end{lstlisting}
